\section{Natural triangular syntax charts}
\label{sec:natural-triangular-charts}

The principal new quartic has $202$ independently nonnegative integer
variables and $154$ quadratic residuals.  Its selected support profile has
$9\,353$ nonzero expanded monomials and coefficient height $1\,854\,167$.
A low-height profile on the same variables has $16\,927$ monomials and
height $416\,488$.  The construction is an exact natural-domain change of
the retained round 21 system; it does not relax natural witnesses to signed
integers or discard a Pell orientation condition.

\begin{center}
\begin{tabular}{lrrrr}
Profile & Variables & Rows & Monomials & Height\\
\hline
Round21 principal &202&154&9\,818&239\,204\,765\\
Round21 low height &202&154&12\,706&31\,145\,294\\
New raw chart &202&154&9\,652&13\,654\,244\\
New principal &202&154&9\,353&1\,854\,167\\
New low height &202&154&16\,927&416\,488\\
New sparse companion &203&155&6\,524&1\,979\,888\\
New sparse companion &204&156&6\,503&1\,981\,520
\end{tabular}
\end{center}

The raw chart already improves both support and height over the preceding
principal profile, and has lower height than its low-height companion.
The last three refinements use independently certified integral residual
shears.  No optimality is asserted.

\subsection{A natural-domain affine chart lemma}

Let $R_1,\ldots,R_r\in\mathbb Z[\boldsymbol x]$, and let $S$ be a set
of gate coordinates admitting an order $g_1,\ldots,g_s$.  Suppose the
retained residuals include
\begin{equation}
 x_{g_i}-f_i(\boldsymbol x)=0,\qquad 1\leq i\leq s,
 \label{eq:triangular-gate-definitions}
\end{equation}
where every coefficient of $f_i$ is nonnegative, $f_i$ does not contain
$x_{g_i}$, and every coordinate from $S$ occurring in $f_i$ belongs to
$g_1,\ldots,g_{i-1}$.  Coordinates outside $S$ require no new interpretation.

Perform elementary changes, read sequentially,
\begin{equation}
 x_v^{\rm new}=x_v^{\rm current}+L(\boldsymbol x^{\rm current}),
 \label{eq:nonnegative-affine-shear}
\end{equation}
where $v\in S$, the affine polynomial $L$ has nonnegative integer
coefficients and does not contain $x_v$.  Let $F$ be their composition and
$G=F^{-1}$ its integral affine inverse.  Then
\begin{equation}
 F:\{\boldsymbol x\in\mathbb N^n:R_j(\boldsymbol x)=0\ \forall j\}
 \longrightarrow
 \{\boldsymbol y\in\mathbb N^n:R_j(G(\boldsymbol y))=0\ \forall j\}
 \label{eq:natural-chart-bijection}
\end{equation}
is a bijection.

Indeed, every forward step preserves nonnegativity.  Conversely, given a
natural zero $\boldsymbol y$, put $\boldsymbol x=G(\boldsymbol y)\in
\mathbb Z^n $.  Coordinates outside $ S$ are unchanged and nonnegative.
All original residuals vanish by substitution.  Induction in
\eqref{eq:triangular-gate-definitions} now proves $x_{g_i}=f_i(\boldsymbol
x)\geq 0 $.  Both affine composite identities hold over $\mathbb Z^n$, so
these maps are mutually inverse on the natural zero sets.

The inverse is emphatically not required to preserve the whole orthant:
$x_v^{\rm current}=x_v^{\rm new}-L(\boldsymbol x^{\rm new})$ can be
negative at points that are not zeros.  It is the retained triangular
positive gate system that restores nonnegativity.  Since the substitutions
are affine, quadratic residuals remain quadratic.

\subsection{The twelve explicit changes}

Use the inherited formal labels, and abbreviate
$a=\texttt{crt.encoded.a}$, $b=\texttt{crt.encoded.b}$,
$G=\texttt{syntax.shared\_shifted\_square}$,
$m_j=\texttt{mul.}j$, and $a_j=\texttt{add.}j$.  The changes are
\begin{align*}
 m_{36}&\gets m_{36}+2b+3,\\
 m_{97}&\gets m_{97}+8b+m_{36}+4m_{60}+11,\\
 m_{173}&\gets m_{173}+2a+m_{74}+16G+16,\\
 m_{243}&\gets m_{243}+16m_{232}+4m_{237}+33,\\
 m_{131}&\gets m_{131}+2a+4m_{36}+m_{74},\\
 m_{48}&\gets m_{48}+32m_{36}+4G+34,\\
 m_{225}&\gets m_{225}+16m_{214}+4m_{219}+34,\\
 m_{274}&\gets m_{274}+16m_{263}+4m_{268}+34,\\
 m_{148}&\gets m_{148}+16m_{36}+4G+18,\\
 m_{317}&\gets m_{317}+5m_{36},\\
 m_{256}&\gets m_{256}+2a_{253}+2,\\
 a_{253}&\gets a_{253}+a.
\end{align*}
Every right-hand side means its current value at that step.  In particular,
the last change occurs after the change to $m_{256}$.

In the literal parent row numbering, starting at zero, the restoration
order is
\[
\begin{array}{c|rrrrrrrrrrrr}
\text{gate}&m_{36}&m_{48}&m_{97}&m_{131}&m_{148}&m_{173}&m_{225}&m_{243}&a_{253}&m_{256}&m_{274}&m_{317}\\
\text{row}&6&8&14&18&20&23&30&33&35&36&39&45 .
\end{array}
\]
Each listed row is literally a monic output minus a polynomial with
nonnegative coefficients, and translated inputs occur only at earlier
positions of this order.  The finite checker verifies these facts from the
coefficient lists.  No outer semantic coordinate, power denominator, Pell
ordinate, CRT remainder or source value is translated.

For the $203$-coordinate companion omit the absent targets $m_{173}$ and
$m_{131}$.  For the $204$-coordinate companion additionally omit $m_{97}$.
All inputs in the remaining changes still exist; their own defining-row
orders meet the same criterion.

\subsection{A cancellation that explains the improvement}

Let $m$ and $p$ denote the original $m_{36}$ and $m_{48}$.  The original
$m_{55}$ defining polynomial is
\begin{equation}
 (2p+128b+64m+8G+259)(2p+128b+64m+8G+260).
 \label{eq:dense-pairing-example}
\end{equation}
The first change gives $\ell=m+2b+3$, and the sixth gives
\[
 z=p+32\ell+4G+34=p+64b+32m+4G+130.
\]
Consequently \eqref{eq:dense-pairing-example} becomes simply
\begin{equation}
 (2z-1)(2z)=4z^2-2z.
\end{equation}
A dense quadratic row has become a one-variable quadratic without a new
coordinate or a divisibility condition.  Its negative linear coefficient
is harmless: natural restoration uses the original gate polynomial, not
coefficientwise positivity of the transformed row.  The other changes
make analogous affine-accumulator cancellations.

For integral $k$ and distinct $i,j$, an elementary residual change
$R_i\gets R_i+kR_j$ is reversed by subtraction.  Thus every supplied
residual-basis sequence preserves the common zero set.  Over natural inputs
the sum of squares is zero precisely when every residual is zero, so each
newly squared basis has the same natural zeros.  The numerical
sum-of-squares polynomial itself need not remain equal.  The principal,
low-height, $203$- and $204$-variable profiles use respectively
$24,53,25,26$ such reversible row changes.

\subsection{A separate degree-eight dimension tradeoff}

The raw chart also yields a $200$-variable polynomial of exact degree eight
with $152$ residuals, $10\,392$ expanded monomials and height
$13\,654\,244$.  This is not a quartic-dimension improvement.

Write $n=q>0$, $A=\texttt{shared.lower\_A}+\texttt{shared.A\_slack}$,
and $T=2An-n^2-1$.  Retain the three inherited pooled pairs $(Y_i,Z_i)$
and affine output expressions $m_i$, and retain the natural witness $M$.
The three pooled quadratic rows are
\begin{align*}
 R_0&=U-\sum_iY_i^2,\\
 R_1&=K-2\sum_iY_i(m_i+Z_i)-(2A-n)M+nU,\\
 R_2&=\sum_i(m_i+Z_i)^2+M-nK.
\end{align*}
Neither $U$ nor $K$ occurs elsewhere.  Delete them and these three rows,
and retain the single quartic residual
\begin{equation}
 Q=\sum_i(m_i+Z_i-nY_i)^2-TM
   =R_2+nR_1-n^2R_0.
 \label{eq:direct-pooled-quartic}
\end{equation}
Projection takes every natural zero of the quadratic system to a natural
zero of the new system.  Conversely, first restore the syntax gates and
use the unchanged unpooled power authentication to obtain $n>0$, before
interpreting any pooled power.  Set
\[
 U=\sum_iY_i^2,\qquad
 K=2\sum_iY_i(m_i+Z_i)+(2A-n)M-nU.
\]
The first expression is natural and the second is integral.  Equation
\eqref{eq:direct-pooled-quartic} gives
$nK=\sum_i(m_i+Z_i)^2+M\geq0$; hence $K\geq0$.  Thus both eliminated
witnesses are restored in the natural domain.  Integrality of $K$ comes
from its polynomial formula, not an unproved divisibility assertion.

\subsection{Calibration and exact verification}

The gate restoration recovers the entire inherited residual system before
any Pell or semantic theorem is used.  Its strengthened lower frame,
additive frame budgets, and inverse-power argument for the $y+1$ seed are
unchanged.  In particular no deleted orientation is silently re-assumed.
The inherited finite-refutation theorem therefore transfers to every
profile above.  With the fixed effective ZFC presentation and its retained
C222/ZF/ZFC calibration, ZFC proves that each displayed no-natural-zero
sentence is equivalent to $\operatorname{Con}(\mathrm{ZFC})$.

The package contains every residual and every collected integer
coefficient.  An independent standard-library checker uses sparse
exponent maps to verify both affine compositions, every transformed row,
the original positive gate definitions and their dependency order, all
sum-of-squares coefficients, and the elimination identities.  Separate
checks replay and reverse every residual-basis change and reject corrupted
certificates.  A clean rebuild reproduces the coefficient files, and the
printed TeX residuals are parsed back to their exact polynomials.  Normal
and optimized Python executions both pass.  These finite computations
check exact algebra; they are not represented as an end-to-end
proof-assistant formalization of the unbounded inherited theorem.
